\chapter{Requisitos y arquitectura del banco}
\label{chap:arquitectura}

\section{Requisitos}

Los requisitos salen de lo que exige el ensayo de la caracola (capítulo~\ref{chap:modelo}). El eje
debe moverse a velocidad constante y conocida en los dos sentidos, y hay que registrar
simultáneamente, con la misma base de tiempos, la posición, la tensión del cable y el par del motor.
Todo ello de forma segura y repetible.

\subsection{Requisitos funcionales}

\begin{longtable}{@{}L{0.09\textwidth}L{0.85\textwidth}@{}}
\toprule
\textbf{Id.} & \textbf{Requisito} \\
\midrule
\endhead
RF1 & Configurar y habilitar el accionamiento y gobernar su velocidad por Modbus RTU sobre RS-485. \\
RF2 & Leer la velocidad y el par del motor en cada periodo de control. \\
RF3 & Leer la tensión del cable con una célula de carga y un acondicionador digital. \\
RF4 & Detectar los finales de carrera A (arriba) y B (abajo) y detener el movimiento hacia un final activado. \\
RF5 & Ejecutar una rutina de \emph{homing} que localice A, mida el recorrido hasta B y fije el origen en un desplazamiento configurable. \\
RF6 & Ejecutar ciclos de ensayo B$\rightarrow$A$\rightarrow$B: colocarse en el extremo inferior, subir hasta A y volver a bajar, con velocidad y número de repeticiones configurables. \\
RF7 & Enviar al ordenador la telemetría de cada periodo con su marca de tiempo del dispositivo. \\
RF8 & Recibir órdenes y parámetros desde el ordenador con detección de errores. \\
RF9 & Consultar y programar la ganancia y el cero del acondicionador de la célula, sólo con el eje parado. \\
RF10 & Mostrar el estado de la máquina y los datos en tiempo real en un navegador. \\
RF11 & Grabar automáticamente cada ensayo en un fichero con datos brutos y constantes de conversión. \\
RF12 & Visualizar los ensayos grabados y compararlos con el modelo teórico. \\
RF13 & Mantener una única configuración de la geometría del mecanismo, editable desde la interfaz y compartida por el simulador y el visor de ensayos. \\
\bottomrule
\end{longtable}

\subsection{Requisitos no funcionales}

\begin{longtable}{@{}L{0.09\textwidth}L{0.85\textwidth}@{}}
\toprule
\textbf{Id.} & \textbf{Requisito} \\
\midrule
\endhead
RNF1 & \textbf{Determinismo}: periodo de control y muestreo fijo de \SI{50}{\milli\second}, con dispersión inferior a \SI{1}{\milli\second}. \\
RNF2 & \textbf{Latencia de parada}: una orden de parada, un final de carrera o el límite de fuerza deben anular la consigna en el mismo periodo en que se detectan. \\
RNF3 & \textbf{Robustez de la comunicación}: una trama corrupta no debe ejecutar ninguna orden, y el receptor debe resincronizarse solo. \\
RNF4 & \textbf{Recuperación}: la aplicación debe reconectarse sola si se desconecta o reinicia el microcontrolador. \\
RNF5 & \textbf{Reproducibilidad}: cualquier resultado debe poder recalcularse a partir de los ficheros grabados. \\
RNF6 & \textbf{Coste y disponibilidad}: componentes comerciales y herramientas libres. \\
RNF7 & \textbf{Mantenibilidad}: validación de cada subsistema con un programa de prueba propio y código documentado. \\
RNF8 & \textbf{Coherencia del modelo}: una sola implementación de la geometría para todas las páginas, para que el simulador y el análisis no puedan discrepar. \\
\bottomrule
\end{longtable}

\section{Arquitectura general}

La figura~\ref{fig:arquitectura} muestra la arquitectura, organizada en tres niveles:

\begin{enumerate}
  \item \textbf{Nivel de campo}: el servomotor paso a paso con su accionamiento, la transmisión, la
        caracola, el cable y la barra con la masa; la célula de carga intercalada en el cable y los
        dos finales de carrera.
  \item \textbf{Nivel de control}: la placa adaptadora con la Raspberry Pi Pico~2. Ejecuta el lazo
        de \SI{50}{\milli\second} y es el único elemento que actúa sobre la máquina.
  \item \textbf{Nivel de supervisión}: un ordenador con el servidor \fichero{web\_monitor.py}, que
        traduce entre el protocolo binario del puerto USB y mensajes JSON por WebSocket, graba los
        ensayos, guarda la configuración del mecanismo y sirve las tres páginas web.
\end{enumerate}

\begin{figure}[H]
\centering
\begin{tikzpicture}[
  font=\small, >=Stealth,
  blk/.style={draw, rounded corners=2pt, align=center, minimum height=9mm, minimum width=33mm, inner sep=3pt, fill=white},
  bus/.style={<->, thick},
  grp/.style={draw, dashed, rounded corners=4pt, inner sep=6pt}
]
  \node[blk] (nav) at (0,0) {Navegador\\ \scriptsize monitor · ensayos · mecanismo};
  \node[blk] (srv) at (5,0) {Servidor FastAPI\\ \scriptsize web\_monitor.py};
  \node[blk] (csv) at (10,0) {Ensayos y configuración\\ \scriptsize cycles/*.csv · mecanismo.json};
  \draw[bus] (nav) -- node[below,font=\scriptsize,align=center] {WebSocket\\HTTP} (srv);
  \draw[->,thick] (srv) -- (csv);
  \node[grp, fit=(nav)(srv)(csv), label={[font=\scriptsize\bfseries,anchor=south west]north west:Ordenador}] {};
  \node[blk] (click) at (0,-3) {Load Cell 2 Click\\ \scriptsize ZSC31014, dirección 0x28};
  \node[blk, fill=blue!6] (pico) at (5,-3) {Raspberry Pi Pico 2\\ \scriptsize firmware\_posicion\_v2};
  \node[blk] (rs) at (10,-3) {Transceptor RS-485\\ \scriptsize THVD1400};
  \node[blk] (io) at (5,-4.6) {E/S de la placa\\ \scriptsize DI1--4, AIN1--3, DO1--4};
  \draw[bus] (pico) -- node[below,font=\scriptsize] {I\textsuperscript{2}C} (click);
  \draw[bus] (pico) -- node[below,font=\scriptsize,align=center] {UART\\PIO} (rs);
  \draw[bus] (pico) -- (io);
  \node[grp, fit=(click)(pico)(rs)(io), label={[font=\scriptsize\bfseries,anchor=south west]north west:Placa adaptadora}] {};
  \draw[bus] (srv) -- node[right,font=\scriptsize,align=left] {USB CDC\\ tramas + CRC-8} (pico);
  \node[blk] (lc) at (0,-7.4) {Célula de carga\\ \scriptsize puente de galgas};
  \node[blk] (fc) at (5,-7.4) {Finales de carrera\\ \scriptsize A arriba, B abajo (NC)};
  \node[blk] (drv) at (10,-7.4) {Accionamiento\\ \scriptsize servo paso a paso};
  \node[blk] (mot) at (10,-9.0) {Motor + reducción $N$\\ \scriptsize caracola $r(\varphi)$};
  \node[blk] (mec) at (5,-9.0) {Cable, barra\\ \scriptsize y masa de 2 kg};
  \draw[->,thick] (lc) -- node[pos=0.28,right,font=\scriptsize] {señal mV/V} (click);
  \draw[->,thick] (fc) -- (io);
  \draw[bus] (drv) -- node[pos=0.28,right,font=\scriptsize] {Modbus RTU} (rs);
  \draw[->,thick] (drv) -- (mot);
  \draw[->,thick] (mot) -- (mec);
  \draw[dotted,thick] (mec.west) -| (lc.south);
  \node[grp, fit=(lc)(fc)(drv)(mot)(mec), label={[font=\scriptsize\bfseries,anchor=south west]north west:Máquina}] {};
\end{tikzpicture}
\caption{Arquitectura del banco de ensayos.}
\label{fig:arquitectura}
\end{figure}

\section{Flujo de datos y temporización}

El microcontrolador ejecuta un lazo de periodo fijo (sección~\ref{sec:fw-lazo}). En cada vuelta lee
el estado del servo y de la célula, integra la posición, procesa las órdenes recibidas, actualiza
las máquinas de estado de \emph{homing} y de ciclo, aplica las protecciones y escribe la nueva
consigna; además guarda una muestra de la vuelta en un búfer. Cada \SI{200}{\milli\second} (o antes,
si el búfer se llena) el búfer se envía en una sola trama, con una marca de tiempo por muestra. El
servidor desempaqueta las muestras, las reenvía al navegador y, si hay un ensayo en marcha, las
escribe en su fichero. Toda la cadena usa el reloj del microcontrolador, de modo que las gráficas y
los ficheros reflejan los instantes reales de medida y no los de llegada al ordenador.

La geometría del mecanismo sigue el camino contrario: se edita en la página de configuración, se
valida y se guarda en el servidor (\fichero{mecanismo.json}), y desde ahí la leen tanto el simulador
como el visor de ensayos, de modo que el modelo con el que se compara una medida es siempre el
mismo.

\section{Decisiones de diseño}

\begin{longtable}{@{}L{0.19\textwidth}L{0.21\textwidth}L{0.48\textwidth}@{}}
\caption{Decisiones de diseño y alternativas consideradas.}\label{tab:decisiones}\\
\toprule
\textbf{Aspecto} & \textbf{Elección} & \textbf{Justificación / alternativa descartada} \\
\midrule
\endfirsthead
\toprule
\textbf{Aspecto} & \textbf{Elección} & \textbf{Justificación / alternativa descartada} \\
\midrule
\endhead
Microcontrolador & Raspberry Pi Pico 2 & Sustituye al ESP32-S3 inicial. La placa se diseñó para el módulo Pico, y su PIO permite una UART en los pines que imponía el trazado. \\
Modo del servo & Velocidad & Los primeros firmwares usaban el modo par con un regulador sobre la célula. El ensayo cuasiestático y el método de separación del par (sección~\ref{sec:mod-separacion}) exigen velocidad constante e igual en ambos sentidos. \\
Organización del firmware & Un único dueño de cada bus en el núcleo 0 & La versión de doble núcleo dedicaba el núcleo 1 a Modbus con variables compartidas. El lazo único de periodo fijo elimina las condiciones de carrera y hace predecible la temporización. \\
Origen del ciclo & Arranca en B, el extremo inferior & El ensayo tiene sentido físico desde la extensión completa del mecanismo. Como el \emph{homing} deja el eje en A, el ciclo se coloca primero en B, y ese tramo ni cuenta como repetición ni se graba. \\
Extremos del ciclo & Finales de carrera & Reconocer A y B por los finales, y no por la posición integrada, hace que el ciclo funcione esté donde esté el cero y sin depender de la escala de velocidad del accionamiento. \\
Protocolo con el PC & Binario con CRC-8 y telemetría por lotes & Un protocolo de texto ocupa más y es ambiguo de delimitar. Las tramas binarias con cabecera, longitud y CRC se resincronizan solas. \\
Posición & Integración de la velocidad del servo & No requiere cablear el codificador. Su deriva se acota con el \emph{homing} y los finales de carrera (sección~\ref{sec:fw-posicion}). \\
Células de carga & Una sola (ZSC31014) & El NAU7802 se evaluó como segunda célula y se retiró: compartían bus y el ZSC31014 no tolera transacciones ajenas entre la petición y la recogida del dato. \\
Supervisión & Aplicación web local & No requiere licencias, admite varios clientes y separa la operación en vivo del análisis. \\
Modelo del mecanismo & Una implementación única en el servidor & Cada página tenía su copia de las fórmulas y llegaron a discrepar: el monitor seguía con un cilindro mientras el visor usaba la caracola. Ahora la geometría vive en \fichero{mecanismo.json} y las fórmulas en \fichero{mecanismo.js}. \\
Datos grabados & Cuentas brutas + constantes en la cabecera & Permite recalibrar un ensayo sin repetirlo, que es justo lo que ha hecho falta. \\
\bottomrule
\end{longtable}

\section{Seguridad}
\label{sec:arq-seguridad}

El prototipo mueve una masa pequeña a baja velocidad, pero el sistema incluye varias medidas de
seguridad, todas implementadas en el firmware:

\begin{itemize}
  \item \textbf{Finales de carrera normalmente cerrados} conectados entre la entrada y masa, con
        \emph{pull-up} interna. En reposo la entrada lee nivel bajo, y un final pulsado o un cable
        roto dan nivel alto: la rotura de cable se interpreta como final activado (comportamiento a
        prueba de fallos).
  \item \textbf{Parada hacia el final activado}: en cada periodo se anula cualquier consigna que
        empuje hacia un final pulsado, sea cual sea el modo (manual, \emph{homing} o ciclo).
  \item \textbf{Límite de fuerza} configurable: si el valor absoluto de la célula tarada supera el
        umbral, la consigna se anula. Sólo actúa con una lectura válida reciente (menos de
        \SI{150}{\milli\second}).
  \item \textbf{Temporizadores de tramo} de \SI{30}{\second} en el \emph{homing} y en los ciclos.
  \item \textbf{Bloqueo durante la programación de la célula}, que la deja sin medida durante
        cientos de milisegundos: se rechaza si el eje no está parado y, mientras dura, se fuerza la
        consigna a cero en cada vuelta.
  \item \textbf{Orden de parada} que corta también el \emph{homing} y los ciclos, y orden de apagado
        que deshabilita el accionamiento.
\end{itemize}

Estas medidas dependen del software y del enlace Modbus. En una máquina industrial, las normas de
seguridad funcional exigirían una cadena de parada independiente del control \cite{iso13849}, por
ejemplo cableando los finales de carrera a la entrada de desconexión segura del par (STO) del
accionamiento \cite{iec61800}. Se recoge como trabajo futuro en el
capítulo~\ref{chap:conclusiones}.
